% ============================================
% MEM-09 (ZH)
% 中文翻译 — SPEC-935-R356
% ============================================
\section*{\textbf{MEM-09} --- 最初几年}
\markboth{MEM-09 --- 最初几年}{}

\noindent \textit{不说话的孩子}

\noindent \textbf{1920—1930 —— 不说话的孩子}

\vspace{0.5em}\noindent\rule{\textwidth}{0.5pt}\vspace{0.5em}

\noindent 在收容院的最初几年，我不说话。

\noindent 不是不能说。是不值得说。

\noindent 说话意味着回答。回答意味着解释。解释意味着重复那段故事。而那故事——乱葬坑、怪物、黄眼——没有人会理解。

\noindent 于是我一直沉默。

\noindent 护工们喜欢安静的患者。安静的患者不添麻烦。安静的患者不占时间。安静的患者像一件家具——它在那里，占着空间，却什么都不要求。

\noindent 我是一件家具。

\noindent 坐在院子角落里。看着地面。整天。天天如此。

\noindent 怪物在我体内沉睡。偶尔它醒来，伸个懒腰，我便感到胃里一阵刺痛。它在长大。

\noindent \textit{“你要多吃一点，”它说。}

\noindent \textit{“没有食物。”}

\noindent \textit{“有。你只是没看见。”}

\noindent \textit{“在哪里？”}

\noindent \textit{“在别人身上。”}

\noindent 我没有回答。

\noindent 但我明白了。

\vspace{0.5em}\noindent\rule{\textwidth}{0.5pt}\vspace{0.5em}

\noindent 1922年，我杀了瓦格纳。

\noindent 那之后，护工们开始用另一种方式对待我。不靠近。把食物放在牢房门口。避免看我的眼睛。

\noindent 隔离室成了我的家。

\noindent 一间三米乘两米的牢房。一张铁床。一只桶。一扇对着走廊的带栅栏窗户。我待在那里，坐着，在日记里写。

\noindent 时间过去。

\noindent 日子过去。

\noindent 年岁过去。

\noindent 我数不清时间。日记成了我的钟。每翻过一页就是一天。每写满一页就是一个夜晚。

\noindent 怪物说得越来越多。

\noindent \textit{“写。”}

\noindent \textit{“写什么？”}

\noindent \textit{“写一切。写乱葬坑。写牲畜圈。写饥饿。写下来，让某人读到。”}

\noindent \textit{“谁会读？”}

\noindent \textit{“某个人。总有一天。有人会翻开这本日记，然后看见。”}

\noindent \textit{“看见什么？”}

\noindent \textit{“真相。”}

\noindent \textit{“什么真相？”}

\noindent \textit{“没人愿意看的真相。世界是一座乱葬坑。人出生，挨饿，死去，被遗忘。但日记不会遗忘。日记保存。”}

\noindent 于是我写。

\noindent 一页又一页。

\noindent 一年又一年。

\noindent 写到手指流血。

\noindent 写到墨水用尽——然后用血写。

\noindent 日记和我一起长大。

\noindent 它知道一些连我自己都不知道自己知道的事情。

\noindent 有时我翻到某一页旧的，认不出那笔迹。以为是别人的。但那是我的。我的手。我的字。

\noindent 但那些词语……

\noindent 那些词语不是我的。

\noindent 是它的。

\vspace{0.5em}
\noindent\textit{\textbf{↪ 路线:} \rota{MEM-08} \rotasep \rota{MEM-10}}
